\chapter{Splitting Newton into Linear-Solve and Step-Control Policies}

\section{Motivation}

Once local constitutive updates are written as nonlinear problems
\[
  \mathbf{R}(\mathbf{x}) = \mathbf{0},
\]
it is tempting to treat Newton as a single algorithmic block.  That is still
too coarse for an extensible constitutive framework.

In practice, a Newton-like local solve contains at least two secondary
algorithmic choices:
\begin{itemize}
  \item how the linearized correction
        \[
          \mathbf{K}(\mathbf{x}^{(i)}) \Delta \mathbf{x}^{(i)}
          = -\mathbf{R}(\mathbf{x}^{(i)})
        \]
        is solved;
  \item how the candidate step $\Delta\mathbf{x}^{(i)}$ is accepted, damped,
        or line-searched.
\end{itemize}

Those are not merely implementation details.  They are real extension points:
\begin{itemize}
  \item scalar division versus dense LU versus custom direct solvers,
  \item full Newton versus fixed damping versus backtracking,
  \item future PETSc/KSP or PETSc/SNES-based local experiments.
\end{itemize}

Therefore, even when the outer algorithm is ``Newton'', the library should not
hard-code its internal sub-policies.

\section{Architectural Split}

This phase introduces two new compile-time seams in
\texttt{src/materials/local\_problem}.

\subsection{Local Linear Solve Policy}

\texttt{LocalLinearSolver.hh} defines the policy
\[
  \Delta\mathbf{x} = \mathcal{L}(\mathbf{K}, \mathbf{R}),
\]
through the concept \texttt{LocalLinearSolvePolicy}.

The default implementation is:
\begin{itemize}
  \item scalar division for one-dimensional local problems,
  \item dense LU for fixed-size vector problems.
\end{itemize}

This is intentionally simple, but it places the dependency in the right place.
If a future constitutive model wants a custom block elimination or a PETSc/KSP
adapter, the change belongs to this policy layer, not to the constitutive law.

\subsection{Local Step-Control Policy}

\texttt{LocalStepControl.hh} defines the update acceptance layer:
\[
  \Delta\mathbf{x}_{\text{accepted}}
  = \mathcal{S}\bigl(\mathbf{x}, \Delta\mathbf{x}, \|\mathbf{R}\|\bigr).
\]

Two policies are provided:
\begin{itemize}
  \item \texttt{FullStepControl}, which accepts the full Newton correction;
  \item \texttt{BacktrackingResidualStepControl}, which searches for a reduced
        step that decreases the residual norm.
\end{itemize}

The key point is again architectural: line search is not fused into Newton, but
appears as a compile-time policy.

\section{Newton as an Assembler of Policies}

After this change,
\texttt{NewtonLocalSolver<MaxIterations, LinearSolvePolicy, StepControlPolicy>}
is no longer a monolithic Newton implementation.  It is a small orchestrator:
\begin{enumerate}
  \item evaluate residual and Jacobian,
  \item call the injected linear solve policy,
  \item call the injected step-control policy,
  \item update the local unknown.
\end{enumerate}

This is a better match to the long-term goals of the library.  It preserves
static polymorphism in the hot loop while exposing each real algorithmic choice
as a separately replaceable type.

\section{Why This Matters for Continuum Inelasticity at Finite Strains}

Finite-strain constitutive updates are often more fragile than the current
small-strain $J_2$ scalar return:
\begin{itemize}
  \item the local unknown may be vector-valued,
  \item the Jacobian may be poorly conditioned,
  \item unrestricted Newton steps may violate admissibility or internal
        constraints,
  \item different models may benefit from different line-search or trust-region
        strategies.
\end{itemize}

For that reason, treating Newton as an injectable outer policy but hard-coding
its substeps would still be too rigid.  The present split is what keeps the
continuum route genuinely open.

\section{PETSc Compatibility Path}

The local-problem architecture is intentionally broad enough that a future
PETSc-based local solver can fit without changing the constitutive law
interface.  A PETSc adapter would only need to satisfy the same generic solver
contract:
\begin{itemize}
  \item accept a local problem,
  \item evaluate residual/Jacobian through callbacks,
  \item return a solved local state \texttt{solution}.
\end{itemize}

Whether that adapter uses:
\begin{itemize}
  \item a direct dense local solve,
  \item a PETSc \texttt{KSP},
  \item or a PETSc \texttt{SNES} line-search method,
\end{itemize}
is then an algorithmic choice at the solver layer, not a change in
\texttt{PlasticityRelation}, \texttt{Material<>}, or \texttt{ContinuumElement}.

\section{Verification}

The regression suite now checks:
\begin{itemize}
  \item concept conformance for the linear-solve and step-control policies;
  \item convergence of the generic vector local problem under the default
        Newton sub-policies;
  \item convergence of the same local problem under a custom damped step
        control.
\end{itemize}

The last point is important.  It proves that the Newton implementation is not
architecturally special: it already works as a host for interchangeable
sub-policies.

\section{Strategic Outcome}

The constitutive local-solver stack is now factored as:
\begin{itemize}
  \item \texttt{LocalNonlinearProblem}: physics-level local system,
  \item \texttt{LocalLinearSolvePolicy}: linearized correction,
  \item \texttt{LocalStepControlPolicy}: step acceptance,
  \item \texttt{LocalNonlinearSolverPolicy}: outer nonlinear algorithm,
  \item \texttt{ReturnAlgorithm}: constitutive update driver.
\end{itemize}

This is the right level of openness for the next step:
\begin{itemize}
  \item finite-strain local vector problems driven by
        \texttt{ConstitutiveKinematics},
  \item damage/fracture-oriented local updates,
  \item alternative nonlinear solvers and PETSc-based experiments.
\end{itemize}

\section{References}

\begin{thebibliography}{9}

\bibitem{simo1998localsolver}
J.~C. Simo and T.~J.~R. Hughes,
\emph{Computational Inelasticity},
Springer, 1998.

\bibitem{desouza2008localsolver}
E.~A. de Souza Neto, D.~Peri\'{c}, and D.~R.~J. Owen,
\emph{Computational Methods for Plasticity: Theory and Applications},
Wiley, 2008.

\bibitem{kelley1995}
C.~T. Kelley,
\emph{Iterative Methods for Linear and Nonlinear Equations},
SIAM, 1995.

\end{thebibliography}
